\chapter{Introduction}

BikeTools Analyzer is a Streamlit application that reads cycling activities (FIT, GPX, TCX, or straight from Strava,
Wahoo and Polar), stores them in a local SQLite database and turns them into ride analysis, training-load analytics,
physics-based models (aerodynamics, climbing, thermal) and PDF reports. It runs on a rider's own computer, either with
\texttt{streamlit run} or as a packaged desktop program. All data stays on that computer.

The project is meant to be changed by the people who use it. This guide exists so that a change does not need a
long reading session first: it explains the few ideas the whole code base is built on, then gives step-by-step
recipes for the changes people most often want.

\section*{What you can add, and where to look}

\begin{center}
\small
\begin{tabular}{@{}p{5.4cm}p{3.3cm}p{5.2cm}@{}}
\toprule
I want to\ldots & Read & Touches mostly\\
\midrule
add a per-ride number (metric) & \S\ref{rec:metric} & \code{analysis/}, \code{pipeline.py}, \code{db.py}\\
add an analysis to the Training page & \S\ref{rec:training} & \code{analysis/}, \code{viz/training_page.py}\\
add a section to the ride page & \S\ref{rec:ridesection} & \code{viz/ride_analysis.py}\\
add a chart & \S\ref{rec:plot} & \code{viz/plots.py}, \code{training_plots.py}\\
read a new file format & \S\ref{rec:format} & \code{data/}, \code{pipeline.py}\\
connect a new online service & Chapter~\ref{ch:providers} & \code{data/}, \code{viz/sidebar.py}\\
add a profile setting & \S\ref{rec:setting} & \code{db.py}, \code{viz/user_corner.py}\\
add something to the PDF report & \S\ref{rec:report} & \code{report_data.py}, \code{report_pdf.py}\\
change the look & \S\ref{sec:theme} & \code{viz/theme.py}\\
\bottomrule
\end{tabular}
\end{center}

\section*{Three rules that make changes safe}

\begin{enumerate}
\item \textbf{Calculations never touch the screen.} Everything under \code{src/analysis}, \code{src/physics} and
  \code{src/data} is plain Python on pandas objects, with no Streamlit calls. That is what makes it testable and
  reusable by the PDF report. (A few data modules still import Streamlit for caching; see \S\ref{sec:headless}.)
\item \textbf{The original file is the source of truth.} Every ride's file is stored, so any metric can be
  recomputed later. A new metric therefore never needs the rider to re-upload anything.
\item \textbf{Add a test with the change.} The test suite is quick (under a minute) and runs on every push.
  Chapter~\ref{ch:conventions} lists what to test.
\end{enumerate}

\chapter{Getting started}

\section{Run it from source}

You need Python 3.10 or newer (the continuous-integration build uses 3.11).

\begin{lstlisting}[style=plain]
git clone https://github.com/jplanari/BikeToolsAnalyzer
cd BikeToolsAnalyzer
python -m venv .venv
.venv\Scripts\activate          # Windows;  macOS/Linux: source .venv/bin/activate
pip install -r requirements.txt
streamlit run app.py --server.port 8501
\end{lstlisting}

Port 8501 matters: the Strava, Wahoo and Polar sign-in redirects point at \code{http://localhost:8501/}. If another
program holds it, Streamlit silently moves to 8502 and sign-in breaks.

The database file is \texttt{bikethools.db} in the folder you start from (the spelling is historical). Set the
environment variable \code{BIKETOOLS_DB} to keep it elsewhere, for example outside a OneDrive or Dropbox folder,
whose sync clients lock files in ways SQLite dislikes. The desktop build uses its own per-user data folder.

\begin{tip}[Use a scratch database while developing]
\begin{lstlisting}[style=plain]
set BIKETOOLS_DB=C:\temp\dev.db          # Windows (cmd)
$env:BIKETOOLS_DB="C:\temp\dev.db"       # Windows (PowerShell)
export BIKETOOLS_DB=/tmp/dev.db          # macOS / Linux
\end{lstlisting}
Your real rides stay untouched while you experiment.
\end{tip}

\section{Run the tests}

\begin{lstlisting}[style=plain]
pytest                       # everything
pytest tests/test_upload.py  # one file
pytest -k thresholds -q      # by name
\end{lstlisting}

Continuous integration (\code{.github/workflows/python-app.yml}) installs \texttt{requirements.txt}, then runs
\texttt{pytest} on every push. Two kinds of test exist:

\begin{itemize}
\item \textbf{Headless tests} (most of them) import the calculation and data modules directly. If a heavy library is
  missing on your machine, \code{tests/conftest.py} replaces it with a stub, so these tests still run.
\item \textbf{User-interface tests} (\code{tests/test_ui.py}, \code{tests/test_training_page.py}) start the real
  application with Streamlit's \texttt{AppTest}, click through the pages against a seeded database, and fail on any
  exception. They are skipped locally when Streamlit is not installed, but \emph{required} in CI.
\end{itemize}

\section{Make a change}

\begin{enumerate}
\item Create a branch from \texttt{main}: \texttt{git checkout -b my-feature}.
\item Make the change and its tests; run \texttt{pytest}.
\item Push and open a pull request. CI runs the same tests.
\item To try it as the desktop app, run the \emph{Desktop build} workflow (Actions tab) on your branch.
\end{enumerate}

Keep pull requests focused. A new metric, a new chart and a new service are three pull requests, not one; each is
then easy to review and to revert.

\chapter{Architecture}

\section{Layers}

The code is four layers. A module may import from the layers to its right, never from the left.

\begin{center}
\begin{tikzpicture}[
  box/.style={rectangle,rounded corners=3pt,draw=ink!60,fill=#1!12,minimum height=1.25cm,text width=3.1cm,align=center,font=\small},
  arr/.style={-{Stealth[length=2.5mm]},thick,ink!70}]
\node[box=accent] (viz) {\textbf{viz/}\\screens, charts, sidebar};
\node[box=kw,right=0.9cm of viz] (analysis) {\textbf{analysis/ physics/}\\pure calculations};
\node[box=cm,right=0.9cm of analysis] (data) {\textbf{data/}\\files, services, SQLite};
\draw[arr] (viz) -- (analysis);
\draw[arr] (analysis) -- (data);
\draw[arr,bend left=25] (viz.north) to node[above,font=\scriptsize]{reads through \texttt{cached.py}} (data.north);
\end{tikzpicture}
\end{center}

\begin{description}[style=nextline,leftmargin=1.2em]
\item[\code{src/data}] Everything that is read or written: file parsers (\code{fit.py}, \code{gpx.py}, \code{tcx.py}),
  service clients (\code{strava.py}, \code{wahoo.py}, \code{polar.py}), the SQLite layer (\code{db.py}) and the
  \emph{pipeline} (\code{pipeline.py}) that turns a file into a standard table and its summary numbers.
\item[\code{src/analysis}] Pure functions: normalized power and TSS (\code{power.py}), per-ride quality metrics
  (\code{ride_metrics.py}), the performance management chart (\code{training_load.py}), power profile and critical
  power (\code{power_profile.py}), climbs, segments, correlations, workout building.
\item[\code{src/physics}] Models of the rider and the road: aerodynamic drag and CdA (\code{aerodyn.py}), a climbing
  model (\code{climb_model.py}), heat balance (\code{thermal.py}).
\item[\code{src/viz}] Everything with Streamlit or Plotly in it: the page for each mode, the sidebar, charts, the
  look, and the PDF writers.
\end{description}

The entry point is \code{app.py}. It initialises the database, applies the look, finishes any pending sign-in
redirect, draws the sidebar, then draws one of three pages according to the top navigation bar.

\section{Repository map}

\begin{longtable}{@{}p{5.6cm}p{9.3cm}@{}}
\toprule
Path & Role\\
\midrule
\endhead
\code{app.py} & Entry point; chooses the page.\\
\code{src/viz/sidebar.py} & Users, data sources (upload, Strava, Wahoo, Polar), sign-in callbacks.\\
\code{src/viz/ride_analysis.py} & The ride page: parse, summarise, save, then draw sections.\\
\code{src/viz/training_page.py} & The Training page (load, power profile, durability, intensity, deeper).\\
\code{src/viz/user_corner.py} & Profile, thresholds, history repair, Connections form.\\
\code{src/viz/connections.py} & In-app entry of service credentials.\\
\code{src/viz/plots.py}, \code{training_plots.py} & Plotly chart builders.\\
\code{src/viz/figstyle.py} & Keeps chart titles and legends from overlapping (automatic).\\
\code{src/viz/sections.py} & Tab-like selector and \texttt{remembered()} widgets.\\
\code{src/viz/theme.py} & Look (CSS), welcome page, help texts.\\
\code{src/viz/report_pdf.py}, \code{ride_report_pdf.py} & PDF writers (matplotlib).\\
\code{src/data/pipeline.py} & File $\rightarrow$ standard table $\rightarrow$ summary dict.\\
\code{src/data/db.py} & Schema, migrations, every read and write.\\
\code{src/data/cached.py} & Cached reads for the pages.\\
\code{src/data/credentials.py} & Where service client IDs and secrets come from.\\
\code{src/data/*_import.py} & Import workflows per source.\\
\code{src/data/thresholds.py} & FTP and LTHR estimated from data.\\
\code{desktop/} & Launcher and PyInstaller recipe for the packaged app.\\
\code{docs/} & OAuth forwarding page, physics report, this guide.\\
\code{tests/} & One test file per area.\\
\bottomrule
\end{longtable}

\section{How a ride travels through the app}

\begin{center}
\begin{tikzpicture}[
  node distance=0.55cm,
  s/.style={rectangle,rounded corners=3pt,draw=ink!60,fill=codebg,minimum height=0.9cm,text width=11.5cm,align=left,font=\small},
  arr/.style={-{Stealth[length=2mm]},ink!70}]
\node[s] (a) {\textbf{1. Source.} Upload, Strava, Wahoo or Polar delivers file bytes and a file name.};
\node[s,below=of a] (b) {\textbf{2. \code{prepare_ride}} (\code{pipeline.py}). Detect the format from content, parse, clean impossible
  readings, resample to 1\,Hz, add distance, speed and grade. Result: one \texttt{DataFrame}.};
\node[s,below=of b] (c) {\textbf{3. \code{summarize_ride}.} Normalized power, TSS, zones, power curve, decoupling, \ldots\ in one dict.};
\node[s,below=of c] (d) {\textbf{4. \code{db.save_ride}.} Metadata, power curve, zone times, histograms \emph{and the original file} are stored.};
\node[s,below=of d] (e) {\textbf{5. Pages.} Training analytics read the stored numbers, never the files. The ride page works on the DataFrame.};
\foreach \x/\y in {a/b,b/c,c/d,d/e} \draw[arr] (\x) -- (\y);
\end{tikzpicture}
\end{center}

Stored numbers make the Training page fast: it reads rows, it does not re-parse hundreds of files. The ride page,
which needs second-by-second data, works on the DataFrame from step 2.
